% !TeX root = ../../Book1.tex
% 纲要标题：### 7.3 幂零群
% 纲要位置：计划纲要.md，第 314 行。

\section{幂零群}
\label{b1:sec:0703}

% 纲要内容（仅作写作计划，尚非教材正文）：
%
% * 上、下中心列
% * 有限 p 群的幂零性
% * 有限幂零群与 Sylow 子群直积
%

\subsection{上、下中心列}
\label{b1:subsec:0703:01}

可解群允许每一步的商群交换，却不要求该商群与群中其他元素相容。若一层元素不仅彼此交换，而且在模去更低一层以后与整个群交换，就得到更强的条件。

\begin{definition}\label{b1:def:central-series}
群 \(G\) 的\term[xia-zhongxin-lie]{下中心列}[lower central series]定义为
\[
\gamma_1(G)=G,\qquad \gamma_{i+1}(G)=[G,\gamma_i(G)].
\]
其\term[shang-zhongxin-lie]{上中心列}[upper central series]定义为
\[
Z_0(G)=1,\qquad
Z_{i+1}(G)/Z_i(G)=Z\bigl(G/Z_i(G)\bigr).
\]
右式通过商映射的原像定义 \(Z_{i+1}(G)\)。若 \(\gamma_{c+1}(G)=1\)，称 \(G\) 为\term[mi-ling-qun]{幂零群}[nilpotent group]；最小这样的非负整数 \(c\) 是其幂零类。
\end{definition}
\symidx{\(\gamma_i(G)\)：下中心列}
\symidx{\(Z_i(G)\)：上中心列}
这些子群都是特征子群。对下中心列，这由交换子在自同构下的相容性归纳得出；对上中心列，自同构先诱导商群的自同构，再保持该商群的中心。上中心列的定义也可写为
\[
x\in Z_{i+1}(G)\quad\Longleftrightarrow\quad
[g,x]\in Z_i(G)\text{ 对每个 }g\in G\text{ 成立}.
\]

\begin{proposition}\label{b1:prop:central-series-equivalence}
对任意 \(c\ge0\)，以下条件等价：
\begin{enumerate}
\item \(\gamma_{c+1}(G)=1\)；
\item \(Z_c(G)=G\)；
\item 存在正规子群链 \(1=C_0\le C_1\le\cdots\le C_c=G\)，使 \([G,C_i]\le C_{i-1}\)。
\end{enumerate}
每个幂零群都可解。
\end{proposition}
\begin{proof}
第一条件成立时，取 \(C_i=\gamma_{c+1-i}(G)\)，即得第三条件。若第三条件成立，由 \(C_0=Z_0(G)\) 和上述中心的判别逐步得到 \(C_i\le Z_i(G)\)，所以第二条件成立。第二条件成立时，直接取 \(C_i=Z_i(G)\) 得到第三条件；由 \(\gamma_1(G)=C_c\) 及交换子的单调性可归纳推出 \(\gamma_j(G)\le C_{c+1-j}\)，于是得到第一条件。

第三条件的每个因子 \(C_i/C_{i-1}\) 都阿贝尔，因此 \(G\) 可解。
\end{proof}
非平凡阿贝尔群的幂零类为一。四元数群满足
\[
Q_8'=Z(Q_8)=\{1,-1\},\qquad [Q_8,\{1,-1\}]=1,
\]
所以幂零类为二。相反，\(S_3\) 虽然可解，中心却平凡；其上中心列永远停留在单位群，故不幂零。由交换子的包含关系还可得幂零性传给子群及商群，但“可解扩张仍可解”不能换成“幂零扩张仍幂零”：\(C_3\trianglelefteq S_3\) 的核与商都阿贝尔就是反例。

\subsection{有限 \texorpdfstring{\(p\)}{p} 群的幂零性}
\label{b1:subsec:0703:02}

\begin{theorem}\label{b1:thm:p-group-nilpotent}
每个有限 \(p\) 群都幂零。若群的阶为 \(p^a\)，则幂零类至多为 \(a\)。
\end{theorem}
\begin{proof}
\chapref{b1:ch:04}的类方程已经说明，非平凡有限 \(p\) 群的中心非平凡。如果 \(Z_i(G)\ne G\)，那么 \(G/Z_i(G)\) 仍是非平凡有限 \(p\) 群，故其中心非平凡，得到严格包含 \(Z_i(G)<Z_{i+1}(G)\)。每次严格增大至少使阶乘上 \(p\)；从一阶增大到 \(p^a\) 至多需要 \(a\) 步。因此 \(Z_a(G)=G\)。
\end{proof}
这里并非只使用“中心非平凡”一次，而是在每一层商群中重新应用它。一般非平凡中心并不足以保证幂零，例如 \(C_2\times S_3\) 的中心非平凡，但其商群含有不可幂零的 \(S_3\)。

\ProseParagraphBreak
对阶为 \(8\) 的二面体群 \(D_8=\langle r,s\mid r^4=s^2=1,\ srs^{-1}=r^{-1}\rangle\)，有
\[
[s,r]=r^{-2}=r^2,\quad
\gamma_2(D_8)=\langle r^2\rangle,\quad
\gamma_3(D_8)=1.
\]
它与 \(Q_8\) 同为二类幂零群，但两群不相互同构：前者有五个二阶元，后者只有一个。幂零类衡量的是交换子消失的层数，不是群的完全分类资料。

\subsection{有限幂零群与 Sylow 子群直积}
\label{b1:subsec:0703:03}

有限群的幂零性还可以从 Sylow 子群是否彼此独立看出来。证明先用一个对无限幂零群也成立的观察。

\begin{lemma}\label{b1:lem:nilpotent-normalizer}
幂零群的每个真子群 \(H<G\) 都严格包含在其正规化子 \(N_G(H)\) 中。
\end{lemma}
\begin{proof}
取最小的 \(i\) 使 \(Z_i(G)\not\le H\)，再取 \(z\in Z_i(G)\setminus H\)。此时 \(Z_{i-1}(G)\le H\)。对任意 \(h\in H\)，有 \([z,h]\in Z_{i-1}(G)\le H\)，于是 \(zhz^{-1}\in H\)。对 \(z^{-1}\in Z_i(G)\) 作同样论证，得到 \(zHz^{-1}=H\)。所以 \(z\in N_G(H)\setminus H\)。
\end{proof}

\begin{theorem}\label{b1:thm:finite-nilpotent-sylow}
对有限群 \(G\)，以下条件等价：\(G\) 幂零；\(G\) 的每个 Sylow 子群都正规；\(G\) 是其各个 Sylow 子群的内直积。
\end{theorem}
\begin{proof}
设 \(G\) 幂零，\(P\) 是一个 Sylow 子群，令 \(N=N_G(P)\)。在 \(N\) 中，\(P\) 是正规 Sylow 子群，故是唯一的对应素数的 Sylow 子群，从而是 \(N\) 的特征子群。任何正规化 \(N\) 的元素都保持 \(P\)，所以
\(N_G(N)\le N_G(P)=N\)。若 \(N<G\)，这与上一引理矛盾。因此 \(N=G\)，\(P\trianglelefteq G\)。

反过来，设各 Sylow 子群 \(P_1,\ldots,P_t\) 正规。对 \(i\ne j\)，正规性给出 \([P_i,P_j]\le P_i\cap P_j=1\)，所以不同因子逐元素交换。乘积映射
\[
P_1\times\cdots\times P_t\longrightarrow G,
\qquad (x_1,\ldots,x_t)\longmapsto x_1\cdots x_t
\]
是同态。其核平凡：若某一 \(x_i\) 等于其余因子乘积的逆，那么其阶一方面是对应素数的幂，另一方面整除其余因子群阶的乘积，只能为一。定义域阶等于 \(|G|\)，故此映射为同构。

最后，有限 \(p\) 群已经证明幂零，而直积的交换子逐分量计算，因此
\(\gamma_k(P_1\times\cdots\times P_t)=\prod_i\gamma_k(P_i)\)。取各因子幂零类的最大值，得到直积幂零。
\end{proof}
例如，阶为 \(12\) 的 \(C_3\times C_4\) 和 \(C_3\times V_4\) 均幂零，\(A_4\) 却不幂零，因为其 Sylow 三子群不唯一。判断有限群时，可先检查 Sylow 子群；计算无限群时，则应回到中心列的定义，不能直接搬用有限 Sylow 直积判据。
